Die Hochschule bietet verschiedene Services zur Beratung und Information an. All diese Stellen sind barrierefrei zu erreichen.

\section{Studienberatung}

Es gibt eine hochschulweite Beratungsstelle, die kompetent berät und gegebenenfalls weiterverweist. Das Beratungsangebot umfasst fachliche Beratung über die einzelnen Studiengänge, Studienfinanzierung und Studienrechtsberatung. Eine psychologische Beratungsstelle und Hilfsangebote für ein behindertengerechtes Studium und Studieren mit Kind sind vorhanden. Ebenso werden Informationen für ausländische Studierende und zum Auslandsstudium zur Verfügung gestellt.

Unter angemessener finanzieller und organisatorischer Unterstützung kann ein Teil der Beratungsverpflichtung an die organisierte Studierendenschaft oder andere Stellen der Hochschule abgetreten werden. Die Effizienz der Beratung ist gewährleistet.

\section{Fachberatung}

Die Fachberatung ist dafür zuständig, dass Studierende ihr Studium zielgerichtet durchführen können. Sie muss insbesondere zu den Themen Studien- und Prüfungsplanung und Studienvertiefung(en)/Spezialisierung(en) sowie zu gängigen Nebenfächern kompetent beraten. Dazu zeigt die Beratungsstelle Beispielstudienpläne auf. Ebenso kann sie über die Anerkennung von Leistungen, die an anderen Hochschulen erbracht wurden und Studienwechseln, sowohl Wechsel des Studiengangs als auch der Hochschule, informieren. Die Beratung findet zeitnah statt.

\section{Betreuung zu Beginn des Studiums}

Alle Erstsemesterstudienden werden vor Studienbeginn mit Hilfe eines Infoheftes und der Möglichkeit einer persönlichen Beratung über ihr zukünftiges Studium informiert. Dieses Infoheft muss nicht analog vorhanden sein, es kann sich auch digital auf einer Website der Hochschule befinden.

Die persönliche Beratung bietet vor allem einen Ausblick über das Studium, informiert über Voraussetzungen und Fristen und leitet bei weiterführenden Fragen an die entsprechenden Beratungsstellen weiter.

Das Infoheft beinhaltet eine Übersicht über die Pflichtveranstaltungen des ersten Jahres, wichtige Ansprechpersonen bzw. Anlaufstellen mit Telefonnummer, E-Mail, Raumnummer und Sprechzeiten. Ebenso sind die wichtigen Termine wie die Rückmeldungsfristen und Prüfungsanmeldungszeiträume angegeben. Ein Lageplan der Hochschule mit Öffnungszeiten der Gebäude befindet sich auch im Infoheft. Im Plan sind insbesondere auch die Bibliothek und die Rechnerräume markiert.

\section{Praktikumsberatung}

Falls in einem Studiengang ein Pflichtpraktikum vorgesehen ist, bietet die Hochschule eine Anlaufstelle für Praktikumsbelange. Diese ermöglicht eine zeitnahe Betreuung der Studierenden. Insbesondere bietet sie regelmäßige Sprechzeiten an. Während den Sprechzeiten ist die Anlaufstelle auch telefonisch oder per Mail zu erreichen.

Die Anlaufstelle dient der Vermittlung von Praktikumsstellen. Sie bietet Informationen und Hilfestellung bei eventuellen Problemen. Dies umfasst Informationen rein organisatorischer Art wie auch praktikumsvorbereitende und praktikumsnachbereitende Informationen und Hilfestellungen.

Desweiteren besitzt die Anlaufstelle Informationen zu eventuellen Fristen, Terminen oder Zulassungsvoraussetzungen und setzt Betroffene frühzeitig darüber in Kenntnis.

\section{Website}

Die Hochschule besitzt eine öffentlich zugängliche, barrierefreie Website. Die für das Studium relevanten Ordnungen sind auf dieser zu finden, inklusive auslaufender Ordnugen oder Studiengänge. Sie sind mindestens in der Sprache des Studiengangs einsehbar.

Die angebotenen Lehrveranstaltungen inklusive Name der Dozierenden, kurzer Beschreibung und etwaiger inhaltlicher Voraussetzungen sind mindestens hochschulöffentlich abrufbar. Eine Liste der Lehrstühle und Lehrpersonen mit angegebener Kontaktmöglichkeit steht öffentlich zur Verfügung.

Auf der Website ist ein Verzeichnis aller vorhandenen Servicestellen vorhanden. Insbesondere wird auf die Fachschaften und deren Webpräsenz verwiesen. Ebenso sind wichtige Termine, wie zum Beispiel Rückmeldefristen, Vorlesungszeiten, Prüfungsanmeldungszeiträume und -termine, Termine von Informationsveranstaltungen etc. zu finden.

Alle Webseiten werden regelmäßig auf ihre Aktualität und Richtigkeit überprüft und entsprechend gepflegt.

\section{Auslandsangebot}

Ein Aufenthalt im Ausland ist in den Studienverlauf integrierbar. Hierbei unterstützt die Hochschule die Studierenden bei der Wahl und dem Kontakt zu einer Gasthochschule. Informationen, insbesondere zu anrechenbaren Veranstaltungen, werden zur Verfügung gestellt.

Studierende, die von einer ausländischen Hochschule kommen, werden bezüglich Visum und anderen Rechtsfragen, Wohnungssuche und Integration unterstützt sowie über geeignete Veranstaltungen beraten.
